\documentclass[a4paper,12pt]{paper}

\usepackage{mycommands}

% Worked-example layout: boxed problem statement followed by a solution.
\newcounter{example}
\renewenvironment{example}[1]{%  % the paper class already defines "example"
  \refstepcounter{example}%
  \begin{mdframed}[linewidth=1pt,innertopmargin=6pt,innerbottommargin=6pt]%
  \noindent\textbf{Example \theexample: #1.}\ }{%
  \end{mdframed}}
\newenvironment{solution}{\par\noindent\textbf{Solution.}\ }{\par\medskip}

\title{Worked examples for Lecture 12: Finite elasticity}
\author{David Al-Attar \\ Michaelmas Term, \the\year}

\begin{document}

\maketitle

\noindent These examples accompany Lecture~12. They are intended as practice in the concrete
calculations that underlie the lecture, and are less demanding than the questions on the
problem sets. Numerical results were checked with the scripts in the \texttt{python}
directory.

\begin{example}{Short descriptive title}
Statement of the problem. Notation follows the lecture notes.
\end{example}

\begin{solution}
Solution text, with displayed equations as in the notes:
\begin{equation}
F_{ij}=\frac{\partial\varphi_{i}}{\partial x_{j}}.
\label{eq:1}
\end{equation}
Reference equations within this file as eq.~(\ref{eq:1}). Results from the lecture are quoted by
name (``the Christoffel equation'') rather than by number.
\end{solution}

\end{document}
